\FrontTop{Abstract}
Large language models can now plan, call tools and act, but in financial analysis a general agent can still invent figures, calculate poorly in words and mistake a failed data source for valid data.
This project designed and built OpenTrading, a harness for financial analysis and market research on Vietnamese equities, as a study in how the program around a model shapes its answers.
The harness runs the agent loop in a terminal and adds rules around the model: a stop gate that rejects answers which claim data that was never fetched, the tracing of figures to tool results, arithmetic run as shell commands, a memory of the investor's constraints that only the user can change, and a permission gate for actions.
A separate data server offers 175 market-data tools whose results state their availability, source and time.
The rules were checked by automated tests, and the harness was compared with two general-purpose harnesses, OpenCode and DeepSeek Harness, on 300 tasks over a frozen data snapshot with the same scoring.
OpenTrading answered 279 tasks correctly (\qty{93.0}{\percent}), against 266 for DeepSeek Harness (\qty{88.7}{\percent}) and 245 for OpenCode (\qty{81.7}{\percent}); on paired tasks an exact sign test gave $p=0.015$ against DeepSeek Harness and $p=3.1\times10^{-7}$ against OpenCode.
Its cost per correct answer was about half that of DeepSeek Harness and less than half that of OpenCode.
The difference came mostly from tasks that depend on rules given in an earlier session, while calculation and retrieval tasks ended within three tasks of each other.
Each task ran once, and the memory tasks were written by the group, so the result is a first measurement rather than a general claim.
